\documentclass[11pt,letterpaper]{article}
\usepackage[T1]{fontenc}
\usepackage{lmodern,amsmath,amssymb,amsthm,mathtools,mathrsfs}
\usepackage[margin=0.95in,headheight=15pt]{geometry}
\usepackage{microtype,booktabs,longtable,array,fancyhdr,xcolor,xurl,enumitem}
\usepackage[colorlinks=true,linkcolor=blue!40!black,citecolor=blue!40!black,urlcolor=blue!40!black]{hyperref}
\usepackage{bookmark}
\numberwithin{equation}{section}
\newtheorem{theorem}{Theorem}[section]
\newtheorem{lemma}[theorem]{Lemma}
\newtheorem{proposition}[theorem]{Proposition}
\newtheorem{corollary}[theorem]{Corollary}
\theoremstyle{definition}\newtheorem{definition}[theorem]{Definition}
\theoremstyle{remark}\newtheorem{remark}[theorem]{Remark}
\newtheorem{example}[theorem]{Example}
\newcommand{\Q}{\mathbb Q}\newcommand{\C}{\mathbb C}\newcommand{\N}{\mathbb N}
\newcommand{\D}{\mathcal D}\newcommand{\NF}{\mathcal N}\newcommand{\HF}{\mathcal H}
\newcommand{\A}{\mathcal A}\newcommand{\CC}{\mathcal C}
\newcommand{\dd}{\,\mathrm d}\newcommand{\ii}{\mathrm i}
\newcommand{\av}[1]{\left\langle #1\right\rangle}
\DeclareMathOperator{\Li}{Li}\DeclareMathOperator{\FP}{FP}
\DeclareMathOperator{\CT}{CT}\DeclareMathOperator{\Res}{Res}
\newcommand{\src}[1]{\nolinkurl{#1}}
\setlength{\emergencystretch}{3em}
\setlength{\parskip}{3pt}
\setlist{itemsep=4pt,topsep=5pt}
\allowdisplaybreaks[2]
\setcounter{tocdepth}{2}
\pagestyle{fancy}\fancyhf{}
\fancyhead[L]{\small A universal beta formula for Gamma normal forms}
\fancyhead[R]{\small ProveIt continuation}
\fancyfoot[C]{\thepage}
\hypersetup{pdftitle={A Universal Beta Formula for Gamma Normal Forms},pdfauthor={Research continuation prepared for Vladimir Reshetnikov with OpenAI assistance},pdfsubject={All functional coefficients, polynomial-background polarization, harmonic moments, Hurwitz and Stieltjes identities}}
\begin{document}
\hypersetup{pageanchor=false}
\begin{titlepage}
\centering
\vspace*{0.7cm}
{\large\sffamily PROVEIT / ANALYSIS / POLYLOGARITHMS\par}
\vspace{1.0cm}
{\Huge\bfseries A Universal Beta Formula\\[0.15em] for Gamma Normal Forms\par}
\vspace{0.7cm}
{\Large All coefficients, exact support and signs,\\[0.2em]
endpoint primitives, and polylogarithmic--Stieltjes transforms\par}
\vspace{0.75cm}
{\large Research continuation prepared for Vladimir Reshetnikov\par}
\vspace{0.2cm}
{Mathematical derivations and computational checks with OpenAI assistance\par}
\vspace{0.2cm}
{11 October 2026\par}
\vspace{0.75cm}
\begin{minipage}{0.94\textwidth}
\small
\textbf{Abstract.}
We give a finite rational formula for every coefficient in the canonical
integration-by-parts normal form of
$(-1)^m\Gamma(x)(\Gamma(x)^{-1})^{(m)}$.
Each coefficient is one polynomial beta integral. The formula determines
its exact support: the previously known factor-parity condition is also
sufficient, and every nonzero coefficient has the sign of the parity of
its highest derivative index. Thus the mixed-square question in the
preceding affine-polarization report is solved for all orders, as are the
higher-multiplicity and coupled sectors of this particular Gamma hierarchy.

The proof transports a constant-coefficient quadratic identity to every
polynomial background by an auxiliary differential operator. The highest
quadratic variation then determines the full normal form, not only its
quadratic sector. A finite entire generating function and an explicit
Euler-homotopy primitive retain all integration endpoints.
The resulting polynomial weights give finite Hurwitz and Lerch endpoint
formulas, generalized-harmonic logarithmic moments, all-order Stieltjes
averages and anchored primitives, and finite polylogarithmic evaluations.
A degree-eight weight yields explicit trigamma, $\gamma_0$, $\gamma_1$,
and log-sine identities.

The classical normal-form complement, the affine square kernel, and the
standard special-function calculus are credited to their sources.
The all-order proofs are separate from 8,350 independent rational
coefficient comparisons through weight 32 and 20 floating-point analytic
diagnostics. Historical priority beyond the inspected sources is not
asserted. No proof-assistant formalization or interval certification is
claimed.
\end{minipage}
\vfill
{\small\sffamily Complete proofs \enspace / \enspace exact rational certificates
\enspace / \enspace additive repository integration\par}
\end{titlepage}
\hypersetup{pageanchor=true}
\pagenumbering{roman}
{\small\setlength{\parskip}{0pt}\tableofcontents}
\clearpage\pagenumbering{arabic}

\section{The continuation problem and the universal formula}
\label{ubg:sec:main}

The immediate predecessor, \emph{Affine Polarization of Gamma Normal
Forms} \cite{Affine}, determines every coefficient of
$u_0^a u_1^k u_r^2$. Its first further-research question asks for the
coupled family $u_0^a u_1^k u_2^b u_3^2$, with $b>0$; its second asks for
nonlinear polarization and higher multiplicities. The earlier
\emph{Exact Harmonic-Tail Reconstruction and Variational Gamma Identities}
\cite{HTG} gives the complete $u_0,u_1,u_2$ projection through a
variational recurrence. We solve these coefficient questions in a
stronger form: every functional monomial of the ordinary inverse-Gamma
hierarchy has an explicit finite formula.

This is a theorem about a specified rational differential-polynomial
normal form. It does not assert that arbitrary individual polygamma
integrals have elementary antiderivatives, or that a set of zeta values
is arithmetically independent. The original Gaussian $S_6$, revised
$S_8$, and general unequal harmonic-tail moment-prefix questions are
not settled here.

Let
\begin{equation}\label{ubg:eq:ring}
 \mathscr R=\Q[u_0,u_1,u_2,\ldots],\qquad
 \D u_j=u_{j+1},\qquad \operatorname{wt}(u_j)=j+1.
\end{equation}
A monomial is \emph{functional} if it is a power of $u_0$, including $1$,
or if its largest positive index occurs at least twice. Write $\NF$
for the unique projection onto the rational span of these monomials,
along total derivatives. Its construction and uniqueness are recalled
in Section~\ref{ubg:sec:normal}. Define
\begin{equation}\label{ubg:eq:hierarchy}
 A_0=1,\qquad A_{m+1}=u_0A_m-\D A_m,\qquad N_m=\NF(A_m).
\end{equation}
Under $u_j=\psi^{(j)}(x)$, with $\psi^{(0)}=\psi$, one has
\begin{equation}\label{ubg:eq:Gamma-hierarchy}
 A_m(\psi(x),\psi'(x),\ldots)
   =(-1)^m\Gamma(x)\frac{\dd^m}{\dd x^m}\frac1{\Gamma(x)}.
\end{equation}

For a polynomial $Q(v)$ put
\begin{equation}\label{ubg:eq:average}
 \av{Q}=\frac12\int_{-1}^1Q(v)\dd v,
 \qquad
 \Theta_j(v)=\frac{(v+1)^{j+1}-(v-1)^{j+1}}
                 {2^{j+1}(j+1)!}\quad(j\ge1).
\end{equation}
The first three polynomials are
\begin{equation}\label{ubg:eq:first-theta}
 \Theta_1=\frac v2,\qquad
 \Theta_2=\frac{1+3v^2}{24},\qquad
 \Theta_3=\frac{v+v^3}{48}.
\end{equation}
All nonzero coefficients of $\Theta_j$ are positive, and
$\Theta_j(-v)=(-1)^j\Theta_j(v)$.

\begin{theorem}[Every functional coefficient]\label{ubg:thm:universal}
Let $d\ge1$, $a,e_1,\ldots,e_d\ge0$, and $e_d\ge2$. Put
\[
 M=u_0^a\prod_{j=1}^d u_j^{e_j},\qquad
 m=a+\sum_{j=1}^d(j+1)e_j,
\]
and define the rational polynomial
\begin{equation}\label{ubg:eq:weight-polynomial}
 P_M(v)=(1-v^2)^d\Theta_d(v)^{e_d-2}
                      \prod_{j=1}^{d-1}\Theta_j(v)^{e_j}.
\end{equation}
Then
\begin{equation}\label{ubg:eq:universal}
 \boxed{\displaystyle
 [M]N_m=
 \frac{(-1)^d m!}
 {2^{2d}d!(d+1)!\,a!\prod_{j=1}^d e_j!}\,\av{P_M}.}
\end{equation}
The coefficient of $u_0^m$ is $1$. A monomial with a unique largest
positive index has coefficient zero. Coefficients at a different total
weight are zero.
\end{theorem}

The average in \eqref{ubg:eq:universal} is a finite rational operation:
if $P_M(v)=\sum q_kv^k$, then
$\av{P_M}=\sum_{k\text{ even}}q_k/(k+1)$. Thus the theorem is not an
implicit recurrence or a numerical integral prescription. It avoids
constructing the full Bell polynomial or reducing all of its partitions.

\begin{corollary}[Exact support and sign]\label{ubg:cor:sign}
For the functional monomial in Theorem~\ref{ubg:thm:universal},
\begin{equation}\label{ubg:eq:support}
 [M]N_m=0\quad\Longleftrightarrow\quad
                      \sum_{j=1}^d je_j\text{ is odd}.
\end{equation}
When the coefficient is nonzero, its sign is exactly $(-1)^d$.
\end{corollary}
\begin{proof}
The integrand has parity $\sum je_j-2d$, which is the parity in
\eqref{ubg:eq:support}. In the even case it is strictly positive for
$0<v<1$, since every $\Theta_j(v)$ is positive there. The prefactor in
\eqref{ubg:eq:universal} has sign $(-1)^d$.
\end{proof}

The parity condition alone was already available from centering
\cite{Manuscript,HTG,Affine}. Its sufficiency and the uniform sign law
follow from the new coefficient formula. In particular, there are no
additional accidental zero coefficients within this functional basis.

\subsection{A positive finite-sum version}
Suppose the parity is even and write
\[
 R_M(v)=\Theta_d(v)^{e_d-2}\prod_{j<d}\Theta_j(v)^{e_j}
                 =\sum_{\ell\ge0}r_\ell v^{2\ell}.
\]
Every $r_\ell\ge0$, and only finitely many are nonzero. The beta integral
\begin{equation}\label{ubg:eq:beta-basic}
 \int_0^1v^{k}(1-v^2)^d\dd v
  =\frac{2^dd!}{\prod_{h=0}^d(k+2h+1)},\qquad \Re k>-1,
\end{equation}
gives the equivalent formula
\begin{equation}\label{ubg:eq:positive-coefficient}
 [M]N_m=
 \frac{(-1)^dm!}{2^d(d+1)!\,a!\prod_j e_j!}
 \sum_{\ell\ge0}\frac{r_\ell}{\prod_{h=0}^d(2\ell+2h+1)}.
\end{equation}
This representation does not subtract alternating coefficients of
$(1-v^2)^d$. It also makes the nonvanishing transparent.

\section{Canonical reduction and the centered Gamma series}
\label{ubg:sec:normal}

The complement used here is classical in differential algebra
\cite[Section 4.4.1]{GRR} and is already established in the ProveIt
manuscript \cite{Manuscript}. We include the elementary proof because
an exact coefficient theorem depends on the precise choice of complement.

\begin{lemma}[Functional decomposition]\label{ubg:lem:normal}
Every $P\in\mathscr R$ has a unique decomposition
\begin{equation}\label{ubg:eq:decomposition}
 P=\NF(P)+\D\HF(P),\qquad \NF(P)\text{ functional},
 \qquad [1]\HF(P)=0.
\end{equation}
Both maps preserve factor degree. On a homogeneous polynomial of
positive weight $m$, the remainder has weight $m$ and the primitive
has weight $m-1$.
\end{lemma}
\begin{proof}
If a monomial has unique largest index $r\ge1$, write it as
$u_r u_{r-1}^aV$, where $V$ uses indices below $r-1$. Then
\begin{equation}\label{ubg:eq:ibp}
 u_r u_{r-1}^aV
 =\D\left(\frac{u_{r-1}^{a+1}V}{a+1}\right)
       -\frac{u_{r-1}^{a+1}\D V}{a+1}.
\end{equation}
Every monomial of the remainder has largest index below $r$.
Repeated replacement terminates. The stated gradings are preserved.
For uniqueness, let a nonconstant $H$ have largest variable index $j$.
The terms $u_{j+1}\partial H/\partial u_j$ in $\D H$ are nonzero,
have largest index $j+1$ exactly once, and cannot cancel terms of
smaller largest index. Hence $\D H$ cannot be functional. The only
remaining freedom in $H$ is a scalar, removed by the normalization.
\end{proof}

The quotient is a vector-space quotient, not an algebra quotient.
For example, $\NF(u_1)=0$ but $\NF(u_1^2)=u_1^2$.
Nor may one truncate high derivatives before reduction:
$\NF(u_0u_4)=u_2^2$. All restrictions to finitely many jet variables
below are performed \emph{after} applying $\NF$.

The exponential generating function of \eqref{ubg:eq:hierarchy} is
\begin{equation}\label{ubg:eq:A-series}
 \A(t)=\sum_{m\ge0}A_m\frac{t^m}{m!}
 =\exp\left(\sum_{j\ge0}\frac{(-1)^j u_jt^{j+1}}{(j+1)!}\right).
\end{equation}
Its centered translate is
\begin{align}
 \CC(t)&=e^{t\D/2}\A(t)
 =\exp\left(\sum_{j\ge0}
       \frac{u_{2j}t^{2j+1}}{2^{2j}(2j+1)!}\right),
 \label{ubg:eq:centered}\\
 \NF\CC(t)&=\NF\A(t)=\sum_{m\ge0}N_m\frac{t^m}{m!}.
 \label{ubg:eq:same-normal}
\end{align}
Indeed, $(e^{t\D/2}-1)P(t)$ is a total derivative coefficientwise.
Under the digamma substitution, \eqref{ubg:eq:A-series} is
$\Gamma(x)/\Gamma(x-t)$, while \eqref{ubg:eq:centered} is
$\Gamma(x+t/2)/\Gamma(x-t/2)$ near $t=0$ away from poles.

Translation $u_0\mapsto u_0+c$ commutes with $\D$ and preserves the
functional span, so it commutes with $\NF$. Consequently
\begin{equation}\label{ubg:eq:Appell}
 \partial_{u_0}\NF\CC(t)=t\NF\CC(t),\qquad
 \partial_{u_0}N_m=mN_{m-1}.
\end{equation}
Let $\sigma_d$ set $u_j=0$ for $j>d$, and put $z_j=u_jt^{j+1}$.
The weight grading and \eqref{ubg:eq:Appell} give unique series $F_d$ with
\begin{equation}\label{ubg:eq:Fd}
 \sigma_d\NF\CC(t)=e^{tu_0}F_d(z_1,\ldots,z_d),
 \qquad F_0=1.
\end{equation}
Normal support implies
\begin{equation}\label{ubg:eq:initial-Fd}
 F_d|_{z_d=0}=F_{d-1},\qquad
 \partial_{z_d}F_d|_{z_d=0}=0.
\end{equation}
The problem is therefore to determine the second derivative of $F_d$
with respect to its highest jet coordinate. This is precisely what
quadratic variation on a polynomial background accomplishes.

\section{Weighted quadratic forms and the affine kernel}
\label{ubg:sec:affine}

We first recover the affine square kernel of \cite{Affine}. This
part is an inherited input, reproduced to fix its normalization and
to expose the transport operation used in the next section.
Let
\begin{equation}\label{ubg:eq:Kt}
 K_t(\xi)=\frac{2\sinh(t\xi/2)}\xi,
 \qquad K_t(0)=t.
\end{equation}
The centered series is formally $\CC(t)=\exp(K_t(\D)u_0)$.
For a perturbation $u_0=q+px+\varepsilon h(x)$, the quadratic
coefficient is
\begin{equation}\label{ubg:eq:quadratic-affine}
 [\varepsilon^2]\CC(t)=\frac12 e^{tq+tpx}(K_t(\partial_x)h)^2.
\end{equation}

\begin{lemma}[Diagonal uniqueness with variable coefficients]
\label{ubg:lem:diagonal}
A finite local quadratic differential expression in $h$ can be reduced,
modulo an ordinary total $x$-derivative, to a sum
$\sum_r a_r(x)(h^{(r)})^2$. The diagonal coefficients are unique.
If a term before reduction is $a(x)h^{(i)}h^{(j)}$, no square of
order greater than $\lfloor(i+j)/2\rfloor$ is needed.
These statements also hold coefficientwise for formal $t$-series
with finite differential order at every coefficient.
\end{lemma}
\begin{proof}
Integration by parts moves a derivative from the larger index to the
smaller one or to the coefficient. The former reduces the difference
of indices, and the latter reduces their sum. Terms of neighboring
orders are reduced by
$2h^{(r)}h^{(r+1)}=((h^{(r)})^2)'$.
This proves existence and the order assertion.

For uniqueness, suppose $\sum_{r=0}^R a_r(x)(h^{(r)})^2$ is a local
total derivative. Its Euler derivative with respect to $h$ is zero.
The coefficient of $h^{(2R)}$ in that Euler derivative is
$2(-1)^R a_R(x)$. Thus $a_R=0$, and descending induction gives all
$a_r=0$. Explicit $x$-dependence of the coefficients causes no problem.
The primitive must be a finite-order local differential expression;
an unrestricted nonlocal antiderivative is not part of this quotient.
\end{proof}

\begin{lemma}[Affine diagonal kernel]\label{ubg:lem:affine}
Let $\kappa$ be an auxiliary constant, independent of $t$. Modulo local
total derivatives,
\begin{equation}\label{ubg:eq:affine-diagonal}
 \frac12e^{\kappa x}(K_t(\partial_x)h)^2
 \equiv e^{\kappa x}\sum_{r\ge0}t^{2r+2}
           g_r(t\kappa/2)(h^{(r)})^2,
\end{equation}
where
\begin{align}
 c_r&=\frac{(-1)^r}{2^{2r+1}r!(r+1)!},
 \label{ubg:eq:cr}\\
 g_r(y)&=c_r\av{(1-v^2)^r e^{yv}}
       =c_r\int_0^1(1-v^2)^r\cosh(yv)\dd v.
 \label{ubg:eq:gr}
\end{align}
Every coefficient of the primitive in \eqref{ubg:eq:affine-diagonal}
can be chosen as $e^{\kappa x}$ times a quadratic differential polynomial
with coefficients in $\Q[\kappa]$.
\end{lemma}
\begin{proof}
Polarize a quadratic expression using derivative symbols $\xi_1,\xi_2$.
A weighted derivative has symbol divisible by
$\xi_1+\xi_2+\kappa$. In the quotient put
$\xi_1=\xi$, $\xi_2=-\xi-\kappa$, and
$\eta=-\xi(\xi+\kappa)$. The symbol of $(h^{(r)})^2$ is $\eta^r$.
Symmetric polynomials subject to this linear relation form the polynomial
ring in $\eta$ over $\Q[\kappa]$.

With $y=t\kappa/2$, the product-to-sum identity gives
\begin{align}
 \frac12K_t(\xi)K_t(-\xi-\kappa)
 &=\frac{\cosh y-\cosh\sqrt{y^2-t^2\eta}}\eta\notag\\
 &=\sum_{r\ge0}\frac{(-1)^rt^{2r+2}}{(r+1)!}
       \left(\partial_{y^2}^{\,r+1}\cosh y\right)\eta^r.
 \label{ubg:eq:affine-symbol}
\end{align}
The entire series $\cosh\sqrt z=\sum_{n\ge0}z^n/(2n)!$ introduces no
branch, and the numerator is divisible by $\eta$.
The beta integral \eqref{ubg:eq:beta-basic}, compared coefficientwise
with the last line, gives exactly \eqref{ubg:eq:gr}. At every power
of $t$, the difference of the two quadratic symbols is divisible by
$\xi_1+\xi_2+\kappa$. Polynomial division gives the asserted local
primitive without introducing denominators in $\kappa$.
\end{proof}

The indices $r=0,1$ in \eqref{ubg:eq:affine-diagonal} are diagonal
\emph{quadratic-variation} coefficients. They are not the constant and
linear coefficients of a normal-form generating function. Confusing
these meanings would destroy the highest-variation argument below.

\section{Polynomial-background transport and the all-order proof}
\label{ubg:sec:transport}

The essential new operation is simple: differential superposition in
the auxiliary affine slope transports the \emph{entire local identity},
including its primitive, to an arbitrary polynomial weight. It is not
a substitution of a variable slope into a constant-coefficient quotient.

\begin{lemma}[Differential superposition]\label{ubg:lem:superposition}
Let $P_t(X)$ be a polynomial in $X$ with formal $t$-series coefficients,
with $P_t(X)=O(t)$. For any constant $a$, coefficientwise in $t$,
\begin{equation}\label{ubg:eq:superposition}
 \left.e^{P_t(\partial_\kappa)}e^{\kappa(x+a)}\right|_{\kappa=0}
                         =e^{P_t(x+a)}.
\end{equation}
The operation commutes with $\partial_x$ and hence transports local
integration-by-parts identities to local integration-by-parts identities.
\end{lemma}
\begin{proof}
The exponential $e^{\kappa(x+a)}$ is an eigenfunction of every polynomial
in $\partial_\kappa$. Expanding the outer exponential proves
\eqref{ubg:eq:superposition}. Because $P_t=O(t)$, only finitely many
powers of that operator are relevant at a specified coefficient of $t$.
The operators $\partial_x$ and $\partial_\kappa$ commute. Locality in
the perturbation and finite differential order at each coefficient are
preserved. The assertion remains valid with $a=tv/2$ and subsequent
integration over the compact parameter interval $-1\le v\le1$.
\end{proof}

Let $b(x)$ be an arbitrary polynomial background, and define its centered
integral
\begin{equation}\label{ubg:eq:background-P}
 P_t(x)=\int_{-t/2}^{t/2}b(x+s)\dd s.
\end{equation}
The integral here means its finite polynomial expansion, so no analytic
assumption on $b$ is necessary. For $u_0=b(x)+\varepsilon h(x)$,
\begin{equation}\label{ubg:eq:quadratic-general}
 [\varepsilon^2]\CC(t)=\frac12e^{P_t(x)}(K_t(\partial_x)h)^2.
\end{equation}
Taylor expansion at $x$ gives the key identity
\begin{equation}\label{ubg:eq:shifted-background}
 P_t(x+tv/2)
 =t b(x)+\sum_{j\ge1}b^{(j)}(x)t^{j+1}\Theta_j(v).
\end{equation}
Indeed, the coefficient of $b^{(j)}(x)$ is
$j!^{-1}\int_{-t/2}^{t/2}(tv/2+s)^j\dd s$, which is
$t^{j+1}\Theta_j(v)$.

\begin{theorem}[Polynomial-background quadratic kernel]
\label{ubg:thm:transport}
If $b$ has degree at most $d$ and $z_j=b^{(j)}(x)t^{j+1}$, the unique
diagonal form of \eqref{ubg:eq:quadratic-general} is
\begin{equation}\label{ubg:eq:general-diagonal}
 e^{tb(x)}\sum_{r\ge0}t^{2r+2}
       G_r(z_1,\ldots,z_d)(h^{(r)})^2,
\end{equation}
where
\begin{equation}\label{ubg:eq:Gr}
 \boxed{\displaystyle
 G_r(\boldsymbol z)=c_r\av{(1-v^2)^r
                    \exp\left(\sum_{j=1}^d z_j\Theta_j(v)\right)}.}
\end{equation}
Each $G_r$ is an entire function of the finitely many $z_j$.
\end{theorem}
\begin{proof}
Apply $e^{P_t(\partial_\kappa)}$ to the affine local identity
\eqref{ubg:eq:affine-diagonal}, and evaluate at $\kappa=0$.
Its left side becomes \eqref{ubg:eq:quadratic-general}.
On the right, write the factor $e^{\kappa x}g_r(t\kappa/2)$ as
\[
 c_r\av{(1-v^2)^r e^{\kappa(x+tv/2)}}.
\]
Lemma~\ref{ubg:lem:superposition} sends this to
$c_r\av{(1-v^2)^r e^{P_t(x+tv/2)}}$.
Now use \eqref{ubg:eq:shifted-background}.
The transported primitive is local, and diagonal uniqueness is supplied
by Lemma~\ref{ubg:lem:diagonal}. Compact integration of an exponential
of a polynomial proves entire dependence on the finite parameter list.
\end{proof}

\begin{proposition}[The highest-variation equation]
\label{ubg:prop:highest}
For every $d\ge1$, the canonical generating function
\eqref{ubg:eq:Fd} satisfies
\begin{equation}\label{ubg:eq:highest}
 \boxed{\displaystyle
 \partial_{z_d}^2F_d(z_1,\ldots,z_d)=2G_d(z_1,\ldots,z_d).}
\end{equation}
\end{proposition}
\begin{proof}
Take a polynomial $b$ of degree at most $d$, but perturb it by an
arbitrary $h$; do not truncate $h^{(j)}$ for $j>d$.
The identity $\CC-\NF\CC=\D H$ implies equality of the diagonal
quadratic variations of the two sides after substitution.

First consider a normal monomial with largest index $r>d$.
Its highest factor must occur at least twice. At order $\varepsilon^2$
it contributes only if it contains exactly two factors above $d$,
both of them $u_r$. Its contribution is already a diagonal square
$(h^{(r)})^2$ of order $r>d$. Monomials with three or more factors
above $d$ do not contribute to quadratic variation at all.

For a monomial involving only $u_0,\ldots,u_d$, the second variation
contains terms $h^{(i)}h^{(j)}$ with $i,j\le d$.
Unless $i=j=d$, Lemma~\ref{ubg:lem:diagonal} shows that its diagonal
reduction has order at most $d-1$. Thus the coefficient of the square
$(h^{(d)})^2$ in the diagonal variation of the normal form is precisely
\[
 \frac12\partial_{u_d}^2\bigl(e^{tu_0}F_d\bigr)
 =\frac12e^{tu_0}t^{2d+2}\partial_{z_d}^2F_d.
\]
Theorem~\ref{ubg:thm:transport} gives the coefficient on the original
centered density as $e^{tu_0}t^{2d+2}G_d$.
Diagonal uniqueness proves the equality after every polynomial-background
substitution. At a fixed $x$, arbitrary values of $b(x),\ldots,b^{(d)}(x)$
are realized by a polynomial of degree $d$. Hence this is an identity
in the independent jet variables, proving \eqref{ubg:eq:highest}.
\end{proof}

\begin{proof}[Proof of Theorem~\ref{ubg:thm:universal}]
By \eqref{ubg:eq:highest}, the coefficient of
$z_1^{e_1}\cdots z_d^{e_d}$ in $F_d$, for $e_d\ge2$, is
\[
 \frac{2c_d}{\prod_{j=1}^d e_j!}
 \av{(1-v^2)^d\Theta_d(v)^{e_d-2}
                         \prod_{j<d}\Theta_j(v)^{e_j}}.
\]
Multiply by the coefficient $1/a!$ from $e^{tu_0}$ and by $m!$ from
the exponential generating convention. Since
$2c_d=(-1)^d/(2^{2d}d!(d+1)!)$, this is \eqref{ubg:eq:universal}.
The initial conditions \eqref{ubg:eq:initial-Fd} supply the remaining
highest-index cases, and $F_0=1$ supplies the pure powers of $u_0$.
\end{proof}

This proof is all-order. The finite rational comparisons reported later
are independent checks of normalizations and implementation, not the
reason the coefficient theorem holds for untested weights.

\section{Entire finite-jet generators and solved coefficient families}
\label{ubg:sec:families}

\subsection{A recursion-free generating function}
Define the entire divided difference
\begin{equation}\label{ubg:eq:E2}
 E_2(z,w)=\frac{e^{zw}-1-zw}{w^2}
       =\sum_{n\ge2}\frac{z^nw^{n-2}}{n!},
 \qquad E_2(z,0)=\frac{z^2}{2}.
\end{equation}
Two integrations of the highest-variation equation, using
\eqref{ubg:eq:initial-Fd}, give the following complete generator.

\begin{corollary}[Finite-jet generating function]\label{ubg:cor:generator}
For every $d\ge0$,
\begin{equation}\label{ubg:eq:generator}
 \boxed{\displaystyle
 F_d=1+\sum_{r=1}^d2c_r\av{(1-v^2)^r
       \exp\left(\sum_{j<r}z_j\Theta_j(v)\right)
       E_2(z_r,\Theta_r(v))}.}
\end{equation}
It is entire in $(z_1,\ldots,z_d)$. All apparent divisions by zeros of
$\Theta_r$ are removable. The infinite sum over $r$ is also meaningful
as a formal $t$-series after $z_j=u_jt^{j+1}$, because the $r$th
summand starts at weight $2r+2$.
\end{corollary}

This formula closes the full finite-jet coefficient problem. In
particular, the higher powers of the highest derivative are determined
by integrating its quadratic variation twice in that coordinate.
There is no need to compute a cubic or quartic symbol in several
independent derivative variables. Such symbols remain useful for other
hierarchies, but they are not an unresolved step for the coefficients
in \eqref{ubg:eq:hierarchy}.

For $d=1$, write $y=z_1/2$. Equation~\eqref{ubg:eq:generator} recovers
the predecessor's low sector,
\begin{equation}\label{ubg:eq:F1}
 F_1(2y)=\frac12\left(\cosh y+\frac{\sinh y}{y}
                      -y\int_0^y\frac{\sinh w}{w}\dd w\right),
\end{equation}
with its removable value at the origin. For a quadratic background,
\eqref{ubg:eq:Gr} becomes the particularly simple heat identity
\begin{align}
 G_r(2y,z)
 &=e^{z/24}e^{(z/8)\partial_y^2}g_r(y)\notag\\
 &=c_r\int_0^1(1-v^2)^r e^{z(1+3v^2)/24}\cosh(yv)\dd v.
 \label{ubg:eq:heat}
\end{align}
The scalar factor $e^{z/24}$ is essential: it is the constant term in
the centered integral of the quadratic background. Dropping it would
produce an incorrect mixed coefficient even at weight eleven.

\subsection{The previously open mixed-square family}
\begin{corollary}[All lower quadratic powers with every higher square]
\label{ubg:cor:mixed}
Let $r\ge3$ and $a,k,b\ge0$, and put $m=a+2k+3b+2r+2$.
For even $k$,
\begin{equation}\label{ubg:eq:mixed}
 \boxed{\displaystyle
 [u_0^a u_1^k u_2^b u_r^2]N_m
 =\frac{(-1)^r m!}
 {a!\,2^{k+r+1}k!(r+1)!\,b!24^b}
 \sum_{j=0}^b\frac{\binom bj3^j}
                 {\prod_{\ell=0}^r(k+2j+2\ell+1)}.}
\end{equation}
For odd $k$ the coefficient is zero.
\end{corollary}
\begin{proof}
In \eqref{ubg:eq:universal}, set $e_r=2$, $e_1=k$, $e_2=b$, and
all other exponents zero. Expand $(1+3v^2)^b$ and apply
\eqref{ubg:eq:beta-basic}. The factor $1/2!$ for the highest square
cancels the factor $2$ in the highest-variation equation.
\end{proof}

At $b=0$ this is the affine theorem of \cite{Affine}.
At $r=3$ and every $b>0$ it answers that report's first explicit
coefficient question. The restriction $r\ge3$ is important: when
$r=2$, the two displayed factors involving $u_2$ must be combined,
changing their factorial normalization.
The complete combined $u_2$ sector is instead the following formula.

\begin{corollary}[A finite replacement for the three-variable recurrence]
\label{ubg:cor:u2}
Let $b\ge2$, $a,k\ge0$, and $m=a+2k+3b$. For even $k$,
\begin{equation}\label{ubg:eq:u2}
 [u_0^a u_1^k u_2^b]N_m
 =\frac{m!}{a!\,2^kk!\,b!24^{b-1}}
 \sum_{j=0}^{b-2}
 \frac{\binom{b-2}j3^j}
 {(k+2j+1)(k+2j+3)(k+2j+5)}.
\end{equation}
Odd $k$ gives zero. This evaluates the all-$u_2$ family previously
specified by the variational recurrence in \cite[Section 4]{HTG}.
\end{corollary}
\begin{proof}
Use $d=2$ in Theorem~\ref{ubg:thm:universal} and the same beta integral.
\end{proof}

\subsection{Nonlinear sectors and explicit normal forms}
The universal theorem is not confined to highest squares. For example,
\begin{equation}\label{ubg:eq:cubic}
 \boxed{[u_1u_3^3]N_{14}=-\frac{3185}{3}.}
\end{equation}
Here $\Theta_1\Theta_3=(v^2+v^4)/96$, so the calculation is a sum of
two elementary beta integrals. In particular, the normal form is not
in general integral over $\mathbb Z$, even though the original Bell
polynomial $A_m$ has integer coefficients.
The first mixed-square values include
\begin{equation}\label{ubg:eq:small-mixed}
 [u_2u_3^2]N_{11}=-55,\qquad
 [u_2^2u_3^2]N_{14}=-3640,\qquad
 [u_2u_4^2]N_{13}=91.
\end{equation}
The entire low normal forms are useful normalization checks:
\begin{align}
 N_4={}&u_0^4-u_1^2,\label{ubg:eq:N4}\\
 N_8={}&u_0^8-70u_0^4u_1^2-7u_1^4
                         +28u_0^2u_2^2-u_3^2,\label{ubg:eq:N8}\\
 N_{11}={}&u_0^{11}-330u_0^7u_1^2-1155u_0^3u_1^4
             +462u_0^5u_2^2+990u_0u_1^2u_2^2\notag\\
 &+550u_0^2u_2^3-165u_0^3u_3^2-55u_2u_3^2+11u_0u_4^2.
 \label{ubg:eq:N11}
\end{align}
These finite polynomials are not presented as historically new by
themselves. The new assertion is the all-monomial formula which
produces every such polynomial without an integration-by-parts search.

\subsection{A support-counting identity}
Let $s_m$ be the number of nonzero monomials of $N_m$.
The exact support theorem gives, as a formal power series,
\begin{equation}\label{ubg:eq:support-generator}
 \sum_{m\ge0}s_mq^m
 =\frac1{1-q}\left\{1+\frac12\sum_{d\ge1}q^{2d+2}
  \left[\prod_{j=1}^d\frac1{1-q^{j+1}}
       +\prod_{j=1}^d\frac1{1-(-1)^jq^{j+1}}\right]\right\}.
\end{equation}
To prove it, reserve two copies of the highest variable $u_d$, allow
arbitrary remaining nonnegative exponents below and at $d$, and apply
the even-parity projector. The factor $(1-q)^{-1}$ accounts for $u_0$.
Each coefficient involves only finitely many $d$.
This counts support, not arithmetic relations among the values obtained
after substituting polygamma functions.

\section{Explicit primitives and Gamma--harmonic endpoint identities}
\label{ubg:sec:primitives}

The normal form is only half of an integration identity. We now give
a finite formula for its retained endpoint primitive.
Let
\begin{equation}\label{ubg:eq:Euler}
 \mathscr E(R)=\sum_{r\ge0}(-\D)^r\frac{\partial R}{\partial u_r}
\end{equation}
be the Euler derivative. All sums are finite for a fixed polynomial.
The commutator
$\partial_{u_r}\D=\D\partial_{u_r}+\partial_{u_{r-1}}$, with the last
term omitted for $r=0$, gives $\mathscr E(\D H)=0$ by cancellation.

\begin{proposition}[An explicit normalized primitive]
\label{ubg:prop:homotopy}
Let $R=\D H$ and decompose $R=\sum_{h\ge1}R_h$ by factor degree.
Then its unique constant-free primitive is
\begin{equation}\label{ubg:eq:homotopy}
 \boxed{\displaystyle
 \HF(R)=\sum_{h\ge1}\frac1h
  \sum_{r\ge1}\sum_{j=0}^{r-1}
   u_{r-1-j}(-\D)^j\frac{\partial R_h}{\partial u_r}.}
\end{equation}
In particular, with $R=A_m-N_m$, the same expression gives the
primitive $H_m$ in $A_m=N_m+\D H_m$.
\end{proposition}
\begin{proof}
Euler's identity for factor degree and repeated integration by parts give
\begin{align*}
 hR_h&=\sum_{r\ge0}u_r\frac{\partial R_h}{\partial u_r}\\
 &=u_0\mathscr E(R_h)+\D\left(
    \sum_{r\ge1}\sum_{j=0}^{r-1}
        u_{r-1-j}(-\D)^j\frac{\partial R_h}{\partial u_r}\right).
\end{align*}
Because $\D$ preserves factor degree, each $R_h$ is itself a total
derivative and its Euler derivative vanishes. Division by $h$ proves
the formula. Every displayed summand has positive factor degree, so
there is no scalar term; uniqueness follows from
Lemma~\ref{ubg:lem:normal}.
\end{proof}

For $0<\alpha<\beta$, substitution in the exact polynomial identity yields
\begin{align}
 &\int_\alpha^\beta(-1)^m\Gamma(x)
             \frac{\dd^m}{\dd x^m}\frac1{\Gamma(x)}\dd x\notag\\
 &\hspace{1cm}=\int_\alpha^\beta N_m(\psi(x),\psi'(x),\ldots)\dd x
      +H_m(\psi(\beta),\ldots)-H_m(\psi(\alpha),\ldots).
 \label{ubg:eq:ordinary-identity}
\end{align}
All integrals here are ordinary convergent integrals.
For instance $H_4=-2u_0^3+4u_0u_1-u_2$, so
\begin{equation}\label{ubg:eq:fourth-identity}
 \int_\alpha^\beta
 \left\{\Gamma(x)(\Gamma(x)^{-1})^{(4)}-\psi(x)^4+\psi'(x)^2\right\}\dd x
 =\left[-2\psi^3+4\psi\psi'-\psi''\right]_\alpha^\beta.
\end{equation}
The supplied rational JSON certificates give $A_m,N_m,H_m$ explicitly
for weights $4,8,11,14,18$; the algorithms work at every finite weight.

At positive integer endpoints the entries are harmonic-tail polynomials:
\begin{align}
 \psi(N+1)&=H_N-\gamma,\notag\\
 \psi^{(j)}(N+1)&=(-1)^{j+1}j!
          \bigl(\zeta(j+1)-H_N^{(j+1)}\bigr),\qquad j\ge1,
 \label{ubg:eq:harmonic-endpoints}
\end{align}
where $H_N^{(q)}=\sum_{n=1}^Nn^{-q}$. These are the standard
polygamma recurrence and zeta specialization \cite{DLMFpsi}.
Equations~\eqref{ubg:eq:universal}, \eqref{ubg:eq:homotopy}, and
\eqref{ubg:eq:harmonic-endpoints} therefore specify all endpoint constants
in \eqref{ubg:eq:ordinary-identity} by finite rational operations in
harmonic numbers and zeta values.

\paragraph{Fixed-coordinate finite parts.}
If the lower endpoint is zero, retain the coordinate $x$ and define the
finite part by the constant term of the cutoff expansion after removing
negative powers and logarithmic divergences. Since $H_m(\psi(x),\ldots)$
has a Laurent expansion at zero, the corresponding identity is
\begin{equation}\label{ubg:eq:FP-endpoint}
 \FP\int_0^\beta(A_m-N_m)(\psi(x),\ldots)\dd x
  =H_m(\psi(\beta),\ldots)
                     -[x^0]H_m(\psi(x),\ldots).
\end{equation}
The last term is generally not zero. This is the same endpoint convention
as in the manuscript and predecessor \cite{Manuscript,Affine}; no
coordinate-invariant finite-part assertion is intended.

\paragraph{What this does not imply.}
The theorem evaluates a particular differential-polynomial combination
and its primitive. It does not turn each remaining normal monomial into
a total derivative. For example, functional monomials were chosen
precisely as a complement to total derivatives. Claims about individual
integrals require additional analytic identities, such as the transforms
in the next sections.

\section{Beta logarithmic moments and generalized harmonic numbers}
\label{ubg:sec:logmoments}

The quadratic-background coefficients naturally introduce the weights
\begin{equation}\label{ubg:eq:Prb}
 P_{r,b}(v)=(1-v^2)^r(1+3v^2)^b,
                  \qquad r,b\in\mathbb Z_{\ge0}.
\end{equation}
Their ordinary logarithmic moments have a finite all-order expression.
For a finite alphabet $\boldsymbol x=(x_0,\ldots,x_r)$, let
$h_m(\boldsymbol x)$ be the complete homogeneous polynomial defined by
\begin{equation}\label{ubg:eq:complete-homogeneous}
 \prod_{\ell=0}^r(1-x_\ell t)^{-1}
          =\sum_{m\ge0}h_m(\boldsymbol x)t^m.
\end{equation}
This notation includes $h_0=1$.

\begin{proposition}[Every logarithmic beta moment]
\label{ubg:prop:beta-logs}
For $\Re k>-1$, $m\ge0$, and $r,b\ge0$,
\begin{align}
 &\int_0^1v^kP_{r,b}(v)(-\log v)^m\dd v\notag\\
 &=m!\,2^rr!\sum_{j=0}^b\binom bj3^j
 \frac{h_m\bigl((k+2j+1)^{-1},\ldots,(k+2j+2r+1)^{-1}\bigr)}
      {\prod_{\ell=0}^r(k+2j+2\ell+1)}.
 \label{ubg:eq:beta-logs}
\end{align}
The dots in the alphabet mean spacing by two in the denominator.
\end{proposition}
\begin{proof}
Expand $(1+3v^2)^b$ and apply \eqref{ubg:eq:beta-basic}. Repeated
$k$-differentiation is justified uniformly on compact subsets of
$\Re k>-1$ by $v^{\Re k}|\log v|^m$ domination. Differentiating a
product of reciprocal linear factors gives $(-1)^m m!$ times that
product times its complete homogeneous polynomial in the reciprocals.
\end{proof}

The power sums of the alphabet in \eqref{ubg:eq:beta-logs} are polygamma
differences. With $x=(k+2j+1)/2$ and $q\ge1$,
\begin{equation}\label{ubg:eq:polygamma-alphabet}
 \sum_{\ell=0}^r\frac1{(k+2j+2\ell+1)^q}
 =\frac{(-1)^{q-1}}{2^q(q-1)!}
       \left\{\psi^{(q-1)}(x+r+1)-\psi^{(q-1)}(x)\right\}.
\end{equation}
For $k=2n$, $n\ge0$, the same sum is
\begin{align}
 &H_{2(n+j+r+1)}^{(q)}-H_{2(n+j)}^{(q)}\notag\\
 &\hspace{1cm}-2^{-q}
       \left\{H_{n+j+r+1}^{(q)}-H_{n+j}^{(q)}\right\}.
 \label{ubg:eq:harmonic-alphabet}
\end{align}
Together with the Newton identities for $h_m$, these formulas express
all moments in generalized harmonic numbers at integer parameters.
For example $h_1=p_1$, $h_2=(p_1^2+p_2)/2$, and
$h_3=(p_1^3+3p_1p_2+2p_3)/6$, where $p_q$ denotes the power sum.
The analytic statements are classical beta differentiation; their role
here is to turn the new coefficient weights into explicit harmonic and
polygamma identities at every logarithmic order.

\section{Finite Hurwitz endpoint formulas for the coefficient weights}
\label{ubg:sec:Hurwitz}

The universal formula supplies rational polynomial weights, not just
rational coefficients. Their special-function transforms can be
computed by a common finite endpoint calculus. Polynomial-weight
Hurwitz integration is classical; see Espinosa--Moll
\cite{EM1,EM2} and the standard formulas in \cite{DLMFzeta}.
Here we give a self-contained version and apply it to the weights
arising from all Gamma sectors.

Let $P$ be an even polynomial of degree $D$, vanishing at $\pm1$ to
order at least $r\ge1$. Every nonzero weight $P_M$ in
\eqref{ubg:eq:weight-polynomial} has this property with $r=d$.
For real $A>B>0$, define
\begin{equation}\label{ubg:eq:IP}
 I_P(s;A,B)=\int_{-1}^1P(v)\zeta(s,A+Bv)\dd v,
 \qquad \mu_P=\int_{-1}^1P(v)\dd v.
\end{equation}
Unlike the average notation in the coefficient theorem, $I_P$ uses
the \emph{full} integral. We keep this normalization visible in all
residues and examples.

\begin{theorem}[Finite endpoints and the Mellin bridge]
\label{ubg:thm:endpoints}
With $d_j=P^{(j)}(1)$,
\begin{equation}\label{ubg:eq:endpoints}
 \boxed{\displaystyle
 I_P(s;A,B)=\sum_{j=r}^D
 \frac{(-1)^j d_j}
      {B^{j+1}(1-s)_{j+1}}
 \left\{\zeta(s-j-1,A+B)-(-1)^j\zeta(s-j-1,A-B)\right\}.}
\end{equation}
Here $(a)_n=a(a+1)\cdots(a+n-1)$ is the rising factorial.
The equality is meromorphic in $s$. The only possible pole of the
whole expression is $s=1$, with residue $\mu_P$.
For $\Re s>1$ it also equals the ordinary convergent Mellin integral
\begin{equation}\label{ubg:eq:Mellin}
 I_P(s;A,B)=\frac1{\Gamma(s)}\int_0^\infty
    \frac{t^{s-1}e^{-At}}{1-e^{-t}}
       \left(\int_{-1}^1P(v)e^{-Btv}\dd v\right)\dd t.
\end{equation}
\end{theorem}
\begin{proof}
The standard parameter derivative is
$\partial_a\zeta(s,a)=-s\zeta(s+1,a)$ \cite{DLMFzeta}.
Consequently an antiderivative in $v$ of $\zeta(s,A+Bv)$ is
$\zeta(s-1,A+Bv)/(B(1-s))$.
Repeated integration by parts terminates after degree $D$ and gives
\eqref{ubg:eq:endpoints}, because $P^{(j)}(-1)=(-1)^jd_j$ and
$d_j=0$ for $j<r$.
Initially exclude the finite exceptional set where a displayed
intermediate expression is singular; meromorphic continuation gives
the identity elsewhere.

The Hurwitz zeta function, uniformly for its parameter in the compact
positive interval $[A-B,A+B]$, has only the simple pole at $s=1$ with
residue $1$. Integrating its holomorphic pole-subtracted germ proves
the assertion about the whole sum. This also shows why apparent
singularities of individual endpoint terms must be combined before
evaluation.

Finally insert the standard Mellin representation of $\zeta(s,a)$.
Near zero the absolute integrand is $O(t^{\Re s-2})$; at infinity it
is dominated by a polynomial times $e^{-(A-B)t}$. Fubini's theorem
proves \eqref{ubg:eq:Mellin} in its stated domain.
\end{proof}

At $s\ne1$, the compact integral \eqref{ubg:eq:IP} is an ordinary
integral even when the infinite Mellin integral does not converge.
The latter must not be used outside its convergence domain without an
explicit continuation or subtraction procedure.

For $q\ge1$, the polygamma and Gamma specializations are
\begin{align}
 \int_{-1}^1P(v)\psi^{(q)}(A+Bv)\dd v
       &=(-1)^{q+1}q!\,I_P(q+1;A,B),\label{ubg:eq:psi-transform}\\
 \int_{-1}^1P(v)\log\Gamma(A+Bv)\dd v
       &=\left.\partial_s I_P(s;A,B)\right|_{s=0}
                                      +\frac{\mu_P}{2}\log(2\pi).
 \label{ubg:eq:Gamma-transform}
\end{align}
These follow from
$\psi^{(q)}(a)=(-1)^{q+1}q!\zeta(q+1,a)$ and
$\zeta'(0,a)=\log\Gamma(a)-\tfrac12\log(2\pi)$
\cite{DLMFpsi,DLMFzeta}. Differentiation in $s,A,B$ is legitimate on
compact subsets away from genuine poles and away from the parameter
boundary $A=B$.

\section{Every Stieltjes order and every anchored parameter primitive}
\label{ubg:sec:Stieltjes}

Use the Laurent convention
\begin{equation}\label{ubg:eq:Stieltjes-convention}
 \zeta(1+\varepsilon,a)=\frac1\varepsilon+
      \sum_{m\ge0}\frac{(-1)^m\gamma_m(a)}{m!}\varepsilon^m,
 \qquad \gamma_0(a)=-\psi(a).
\end{equation}
For $j\ge0$, define complete-harmonic coefficients by
\begin{equation}\label{ubg:eq:hj}
 \prod_{\ell=1}^j(1-\varepsilon/\ell)^{-1}
     =\sum_{n\ge0}h_n^{(j)}\varepsilon^n.
\end{equation}
For the empty product, $h_0^{(0)}=1$ and $h_n^{(0)}=0$ for $n>0$.
For example
\[
 h_1^{(j)}=H_j,\qquad
 h_2^{(j)}=\frac{H_j^2+H_j^{(2)}}2,\qquad
 h_3^{(j)}=\frac{H_j^3+3H_jH_j^{(2)}+2H_j^{(3)}}6.
\]
Write $\zeta^{[\ell]}(s,a)=\partial_s^\ell\zeta(s,a)$ for spectral
derivatives, to distinguish them from parameter derivatives.

\begin{theorem}[All-order Stieltjes endpoint formula]
\label{ubg:thm:Stieltjes}
For $P,A,B$ as in Theorem~\ref{ubg:thm:endpoints}, set
\begin{equation}\label{ubg:eq:cj}
 c_j(P,B)=\frac{(-1)^jP^{(j)}(1)}{B^{j+1}j!}.
\end{equation}
Then for every $m\ge0$,
\begin{align}
 &\int_{-1}^1P(v)\gamma_m(A+Bv)\dd v\notag\\
 &=(-1)^{m+1}m!\sum_{j=r}^D c_j(P,B)
    \sum_{\ell=0}^{m+1}\frac{h_{m+1-\ell}^{(j)}}{\ell!}
     \left\{\zeta^{[\ell]}(-j,A+B)
                -(-1)^j\zeta^{[\ell]}(-j,A-B)\right\}.
 \label{ubg:eq:all-Stieltjes}
\end{align}
In particular, the required endpoint spectral derivatives have orders
at most $m+1$, at the finitely many nonpositive integers $-j$.
\end{theorem}
\begin{proof}
At $s=1+\varepsilon$,
\begin{equation}\label{ubg:eq:denominator-Stieltjes}
 \frac1{(1-s)_{j+1}}
   =-\frac1{\varepsilon j!}
                \prod_{\ell=1}^j(1-\varepsilon/\ell)^{-1}.
\end{equation}
Insert this in \eqref{ubg:eq:endpoints}, Taylor-expand the two zeta
values, and compare the coefficient of $\varepsilon^m$ with
\eqref{ubg:eq:Stieltjes-convention}. The finite endpoint sum permits
coefficient extraction term by term. On the compact parameter interval,
the pole-subtracted Hurwitz germ is jointly holomorphic, so its Laurent
coefficients are the ordinary integrals of $\gamma_m$.
\end{proof}

Thus the Stieltjes calculation is an explicit convolution with complete
harmonic numbers, not an unspecified instruction to ``differentiate a
regularized expression.'' The residue identity, obtained at the
previous coefficient, is
\begin{equation}\label{ubg:eq:residue-check}
 -\sum_{j=r}^D c_j(P,B)
       \left\{\zeta(-j,A+B)-(-1)^j\zeta(-j,A-B)\right\}
          =\mu_P.
\end{equation}
It gives a useful exact Bernoulli-polynomial check of all endpoint
normalizations.

\subsection{Arbitrary orders of antiderivation in the center}
For integers $m\ge0$, $d\ge1$, define
\begin{equation}\label{ubg:eq:Stieltjes-primitive}
 \mathcal B_{P;m,d}(A;B)=
 \frac{(-1)^{m+1}m!}{(d-1)!}
 \sum_{\ell=0}^{m+1}\frac{h_{m+1-\ell}^{(d-1)}}{\ell!}
 \left.\partial_s^\ell I_P(s;A,B)\right|_{s=1-d}.
\end{equation}

\begin{theorem}[All anchored primitives]\label{ubg:thm:primitives}
For $A>B>0$,
\begin{equation}\label{ubg:eq:primitive-derivative}
 \partial_A^d\mathcal B_{P;m,d}(A;B)
       =\int_{-1}^1P(v)\gamma_m(A+Bv)\dd v.
\end{equation}
Fix an anchor $A_0>B$. The unique $d$-fold primitive whose first $d$
initial data vanish at $A_0$ is
\begin{equation}\label{ubg:eq:anchored}
 \mathcal B_{P;m,d}(A;B)
 -\sum_{h=0}^{d-1}\frac{(A-A_0)^h}{h!}
            \partial_A^h\mathcal B_{P;m,d}(A_0;B).
\end{equation}
\end{theorem}
\begin{proof}
The repeated Hurwitz parameter derivative gives, meromorphically in $s$,
\[
 \partial_A^d\left\{\frac{I_P(s-d;A,B)}{(1-s)_d}\right\}
                                =I_P(s;A,B).
\]
Set $s=1+\varepsilon$ and use
\[
 \frac1{(-\varepsilon)_d}
 =-\frac1{\varepsilon(d-1)!}
       \prod_{\ell=1}^{d-1}(1-\varepsilon/\ell)^{-1}.
\]
The coefficient of $\varepsilon^m$, multiplied by $(-1)^mm!$, is
\eqref{ubg:eq:Stieltjes-primitive}. Holomorphic dependence on $A$ in
$\Re A>B$ permits commutation of the derivative and coefficient
extraction. Subtracting the Taylor polynomial supplies the stated
initial conditions. The difference of any two such primitives is a
polynomial in $A$ of degree at most $d-1$, hence is zero.
\end{proof}

For $m=0,d=1$, the anchored primitive simplifies to the ordinary Gamma
ratio integral
\begin{equation}\label{ubg:eq:Gamma-anchored}
 -\int_{-1}^1P(v)
       \log\frac{\Gamma(A+Bv)}{\Gamma(A_0+Bv)}\dd v.
\end{equation}
Equations~\eqref{ubg:eq:endpoints} and
\eqref{ubg:eq:Stieltjes-primitive} also give a finite endpoint expression
for every higher primitive. No unchosen polynomial integration constant
is hidden in the notation.

\section{Lerch deformations, polylogarithms, and Fourier endpoint identities}
\label{ubg:sec:polylog}

\subsection{A finite polylogarithmic endpoint calculus}
For $|\xi|<1$, the Lerch series is
\[
 \Phi(\xi,s,a)=\sum_{n\ge0}\frac{\xi^n}{(n+a)^s},\qquad a>0,
\]
where positive real denominators use the real logarithm
\cite{DLMFlerch}. It converges normally in $s$ on every compact set.
Define
\begin{equation}\label{ubg:eq:Lerch-transform}
 L_P(s;A,B;\xi)=\int_{-1}^1P(v)\Phi(\xi,s,A+Bv)\dd v.
\end{equation}
The same finite integration-by-parts proof gives
\begin{equation}\label{ubg:eq:Lerch-endpoints}
 L_P(s;A,B;\xi)=\sum_{j=r}^D
 \frac{(-1)^jd_j}{B^{j+1}(1-s)_{j+1}}
 \left\{\Phi(\xi,s-j-1,A+B)-(-1)^j\Phi(\xi,s-j-1,A-B)\right\}.
\end{equation}
The whole expression is entire in $s$; every apparent exceptional value
is obtained by joint continuation. In $\Re s>0$ its ordinary Mellin
representation replaces $1-e^{-t}$ in \eqref{ubg:eq:Mellin} by
$1-\xi e^{-t}$. The weaker lower-end restriction follows because the
new denominator is nonzero at $t=0$.

If $A\pm B$ are positive integers and $0<|\xi|<1$, the elementary identity
\begin{equation}\label{ubg:eq:integer-Lerch}
 \Phi(\xi,s,N)=\xi^{-N}
 \left\{\Li_s(\xi)-\sum_{n=1}^{N-1}\frac{\xi^n}{n^s}\right\}
\end{equation}
turns \eqref{ubg:eq:Lerch-endpoints} into a finite combination of
polylogarithms at shifted orders and finite harmonic-type sums.
Every spectral derivative and every center primitive follows from the
same identities. The apparent singularity at $\xi=0$ in
\eqref{ubg:eq:integer-Lerch} is removable after the displayed difference
is formed. At $\xi=1$, passage to Hurwitz zeta must respect the genuine
pole and the original convergence domain; the entire-in-$s$ statement
for $|\xi|<1$ is not uniform across that boundary.

\subsection{A finite zeta formula at every complex Fourier order}
The same endpoint polynomial also evaluates a Fourier-polylogarithm
integral without a Hurwitz parameter.

\begin{proposition}[Fourier endpoint identity]\label{ubg:prop:Fourier}
Let $P$ be an even polynomial of degree $D$, with $P(\pm1)=0$.
For $\Re s>1$,
\begin{equation}\label{ubg:eq:Fourier-polylog}
 \boxed{\displaystyle
 \int_{-1}^1P(v)\Li_s(e^{2\pi\ii v})\dd v
 =2\sum_{\ell=1}^{D/2}
      \frac{(-1)^{\ell+1}P^{(2\ell-1)}(1)}{(2\pi)^{2\ell}}
           \zeta(s+2\ell).}
\end{equation}
The integer upper limit is harmless because $D$ is even; zero
endpoint derivatives may be omitted. The identity may be differentiated
in $s$ any fixed number of times within this half-plane.
\end{proposition}
\begin{proof}
For $n\ge1$, repeated integration by parts gives the exact terminating
Fourier coefficient
\[
 \int_{-1}^1P(v)e^{2\pi\ii nv}\dd v
 =2\sum_{\ell=1}^{D/2}
       \frac{(-1)^{\ell+1}P^{(2\ell-1)}(1)}{(2\pi n)^{2\ell}}.
\]
The even endpoint derivatives cancel between $1$ and $-1$, and the
last constant derivative has zero integral over the integer number
of periods. Now use the absolutely uniformly convergent unit-circle
series for $\Li_s$ when $\Re s>1$ \cite{DLMFpolylog}. Summing $n^{-s}$
times the finite Fourier coefficient gives
\eqref{ubg:eq:Fourier-polylog}. Spectral derivatives introduce only
powers of $\log n$, preserving normal convergence on smaller compact
subsets of the half-plane.
\end{proof}

At $s=1$, the ordinary integrable logarithmic boundary value gives
\begin{equation}\label{ubg:eq:logsine-general}
 \int_{-1}^1P(v)\log(2|\sin\pi v|)\dd v
 =-2\sum_{\ell=1}^{D/2}
   \frac{(-1)^{\ell+1}P^{(2\ell-1)}(1)}{(2\pi)^{2\ell}}
                           \zeta(2\ell+1).
\end{equation}
One rigorous route is to insert a radial factor $0<\rho<1$ into the
Fourier series and use $\Li_1(z)=-\log(1-z)$.
The real part converges in $L^1$ to $-\log(2|\sin\pi v|)$:
near each integer the absolute logarithm is dominated by an integrable
constant plus $|\log|v-n||$, uniformly for $\rho$ sufficiently close
to one. The integrated Fourier series is absolutely summable by its
explicit coefficients. The imaginary part cancels by parity.
No unregulated ordinary-integral statement at $s=0$ or negative orders
is inferred; such boundary values can involve distributional terms.

\section{A degree-eight bridge: four explicit evaluations}
\label{ubg:sec:examples}

The first genuinely mixed-square coefficient $u_2u_3^2$ produces
\begin{equation}\label{ubg:eq:P31}
 P(v)=P_{3,1}(v)=(1-v^2)^3(1+3v^2)
                 =1-6v^4+8v^6-3v^8.
\end{equation}
On the interval $1\le x\le3$ put
\begin{equation}\label{ubg:eq:W}
 W(x)=P(x-2)=(x-1)^3(3-x)^3\bigl(1+3(x-2)^2\bigr).
\end{equation}
Its mass is $\int_1^3W(x)\dd x=128/105$.
The endpoint derivatives provide a compact exact certificate:
\begin{center}
\small
\begin{tabular}{c r r r r r r}
\toprule
$j$&3&4&5&6&7&8\\
\midrule
$P^{(j)}(1)$&$-192$&$-2304$&$-14400$&$-54720$&$-120960$&$-120960$\\
$(-1)^jP^{(j)}(1)/j!$&$32$&$-96$&$120$&$-76$&$24$&$-3$\\
\bottomrule
\end{tabular}
\end{center}
Here $A=2$, $B=1$ in the compact Hurwitz transform. Integer translation
reduces its finite endpoint formula to
\begin{equation}\label{ubg:eq:I31}
 I_P(s;2,1)=\sum_{j=3}^8\frac{(-1)^jP^{(j)}(1)}{(1-s)_{j+1}}
   \left\{(1-(-1)^j)\zeta(s-j-1)-1-2^{j+1-s}\right\}.
\end{equation}
This one meromorphic identity proves the following evaluations.

\subsection{Trigamma and the zeroth Stieltjes constant}
\begin{align}
 \boxed{\displaystyle
 \int_1^3W(x)\psi'(x)\dd x
  =-\frac{32}{35}-\frac{48}{\pi^2}\zeta(3)
          +\frac{900}{\pi^4}\zeta(5)-\frac{1890}{\pi^6}\zeta(7).}
 \label{ubg:eq:trigamma-example}
\end{align}
For completeness, at $s=2+\varepsilon$,
\[
 \frac1{(1-s)_{j+1}}
  =\frac1{\varepsilon(j-1)!}
       \left\{1+(H_{j-1}-1)\varepsilon+O(\varepsilon^2)\right\}.
\]
The table shows that the total residue in \eqref{ubg:eq:I31} is zero.
Its elementary finite part is $-32/35$; all $\log2$ terms cancel.
For odd $j=2n+1$, the zeta derivative is reduced by the standard
functional-equation consequence \cite{DLMFintegers}
\[
 \zeta'(-2n)=\frac{(-1)^n(2n)!}{2(2\pi)^{2n}}\zeta(2n+1).
\]
This gives \eqref{ubg:eq:trigamma-example} with no appeal to a numerical
integer-relation guess.

The constant term at the genuine pole $s=1$ is
\begin{equation}\label{ubg:eq:gamma0-example}
 \boxed{\displaystyle
 \int_1^3W(x)\gamma_0(x)\dd x
 =-\frac{547}{1800}-64\zeta'(-3)-240\zeta'(-5)-48\zeta'(-7).}
\end{equation}
Since $\gamma_0=-\psi$, this is also an explicit weighted digamma
integral. The derivatives in this and the next display are derivatives
of the ordinary Riemann zeta function in its argument.

\subsection{An explicit first-Stieltjes integral}
The next coefficient of \eqref{ubg:eq:I31} gives
\begin{align}
 \int_1^3W(x)\gamma_1(x)\dd x
 ={}&\frac{3499687}{10584000}-\frac{32}{21}\log2
       +\frac{352}{3}\zeta'(-3)+548\zeta'(-5)
                       +\frac{4356}{35}\zeta'(-7)\notag\\
   &+32\zeta''(-3)+120\zeta''(-5)+24\zeta''(-7).
 \label{ubg:eq:gamma1-example}
\end{align}
To check all terms without numerical approximation, put
$c_j=(-1)^jP^{(j)}(1)/j!$ and expand
\[
 B_j(\varepsilon)=(1-(-1)^j)\zeta(-j+\varepsilon)
                               -1-2^je^{-\varepsilon\log2}.
\]
Then $I_P(1+\varepsilon;2,1)$ is
\[
 -\frac1\varepsilon\sum_{j=3}^8c_j B_j(\varepsilon)
             \left\{1+H_j\varepsilon
                  +\tfrac12(H_j^2+H_j^{(2)})\varepsilon^2+\cdots\right\}.
\]
The coefficient of $\varepsilon$ is the negative of the right side
of \eqref{ubg:eq:gamma1-example}. This proves the displayed rational,
logarithmic, first-derivative, and second-derivative coefficients
simultaneously. Theorem~\ref{ubg:thm:Stieltjes} gives every subsequent
$\gamma_m$ integral by the same finite rule.

\subsection{A function-valued polylogarithm identity and a log-sine value}
For $\Re s>1$, Proposition~\ref{ubg:prop:Fourier} specializes to
\begin{equation}\label{ubg:eq:polylog-example}
 \boxed{\displaystyle
 \int_1^3W(x)\Li_s(e^{2\pi\ii x})\dd x
    =\frac{24}{\pi^4}\zeta(s+4)
      -\frac{450}{\pi^6}\zeta(s+6)
      +\frac{945}{\pi^8}\zeta(s+8).}
\end{equation}
The unit-circle polylogarithm here is defined by its convergent series
or equivalently its radial boundary value from the disk. This is an
identity in a complex order, not only a list of special integer values.
All fixed spectral derivative orders obey the differentiated identity.
The logarithmic endpoint specialization is
\begin{equation}\label{ubg:eq:logsine-example}
 \boxed{\displaystyle
 \int_1^3W(x)\log(2|\sin\pi x|)\dd x
   =-\frac{24}{\pi^4}\zeta(5)
      +\frac{450}{\pi^6}\zeta(7)
      -\frac{945}{\pi^8}\zeta(9).}
\end{equation}
There is an interior logarithmic singularity at $x=2$, but it is
integrable. No finite-part convention is used in any of
\eqref{ubg:eq:trigamma-example}--\eqref{ubg:eq:logsine-example}.
The finite Lerch formula \eqref{ubg:eq:Lerch-endpoints} gives a second,
disk-valued polylogarithmic deformation of the same weight, independently
of the Fourier specialization.

\section{Source audit, corrections, and reproducible verification}
\label{ubg:sec:audit}

\subsection{Provenance and the exact advance}
The research was based on targeted reads of the canonical manuscript
metadata and relevant normal-form material, the harmonic-tail Gamma
continuation \cite{HTG}, and the affine-polarization source
\cite{Affine}. The incoming directory was inspected at the metadata
level. The affine article was retrieved as a full TeX file from the
user's file library; byte identity with a similarly named incoming ZIP
was not established. Other incoming archives were not exhaustively
unpacked. This is a targeted continuation, not a complete audit of the
large and actively changing repository.

The inspected affine TeX file is named
\src{ProveIt_Affine_Gamma_Polarization_2026-10-11.tex}.
Its SHA-256 fingerprint is
\begin{center}\small\ttfamily
36a06cf61ab08597c5526c3ced33af296fb\\
536466974211073769ee41218c751
\end{center}
The file itself records the older repository metadata anchor
\src{92ee3aaf3087c0e0c29b9ed952e777de5d6accd0}.
That is the predecessor's provenance, not a claim that every live
repository read in this session came from that revision.
Observed search results moved through several revisions; the final
additional search results carried
\src{af91e74dfa841331c51e867542345a3643ee0bf0}.
The source manifest in the package records this distinction.

The contributions are separated as follows. The functional complement,
centering, Appell factor, and the affine kernel are inherited and
explicitly proved again for normalization. The polynomial-background
transport, its highest-variation use to determine \emph{all} normal
coefficients, the universal formula, the sufficient support condition,
and the highest-index sign law constitute the present coefficient
advance relative to the inspected sources. The finite $u_2$ sum and
all mixed-square formulas are consequences, not separate conjectural
extensions. The higher-multiplicity sectors are included in the theorem.

Polynomial-weight Hurwitz integration, beta differentiation, and the
Lerch--polylogarithm connection are established methods
\cite{EM1,EM2,DLMFzeta,DLMFlerch,DLMFpolylog}.
The analytic contribution is their explicit application and all-order
normalization for the weights supplied by the universal Gamma formula,
including the worked degree-eight identities. Neither the general
method nor historical priority for each isolated definite integral is
claimed as new. Targeted repository and public-source searches did not
supply the universal formula, but those searches do not establish
exhaustive historical priority.

\subsection{Corrections and safeguards}
\paragraph{A definite bibliographic correction.}
The affine report's bibliography gives pages 456--473 for Guo,
Regensburger, and Rosenkranz \cite{GRR}. The authors' arXiv record gives
\emph{Journal of Pure and Applied Algebra} \textbf{218} (2014),
456--471, with the same DOI. The page range should therefore be changed
to 456--471. This does not affect any mathematical argument.

\paragraph{A repository editorial duplication.}
The initially inspected canonical manuscript README, with blob identity
\src{dc0fa26d0e5c2ede56df37c5640c7319681dfd25}, repeated the paragraph
about the complete 1,032-input first-letter $X$ family, rank 1,008,
and kernel 24. Remove the second copy if that duplication remains in
the integration target. This observation is tied to the inspected blob,
not a claim about every later version of the moving default branch.

\paragraph{No substantive predecessor error established.}
The coefficient formulas used from the two direct predecessors agree
with the present formula in their overlapping sectors. We did not
identify a substantive mathematical error in those inspected formulas.
The following are implementation safeguards, not retrospective
accusations: reduce before truncation; do not multiply projected
normal forms as though the projection were an algebra homomorphism;
keep the $e^{z/24}$ factor in the quadratic-background transport;
combine removable endpoint singularities before evaluation; and retain
the primitive's lower-end constant in fixed-coordinate finite parts.

\subsection{Independent exact checks}
The package contains two different algorithms. The coefficient module
\src{verification/gamma_coefficients.py} uses only polynomial arithmetic
and the universal beta formula. The reference module
\src{verification/independent_ibp.py} generates $A_m$ by
\eqref{ubg:eq:hierarchy} and applies the terminating local reduction
\eqref{ubg:eq:ibp}; it does not call the coefficient formula.

The exact run checks every functional exponent slot through weight 32,
including those predicted to vanish. The original Bell polynomial has
nonzero coefficient at every partition of its weight, so its functional
slots, together with the reference normal form's support, enumerate the
whole required test set. The recorded results are:
\begin{center}
\begin{tabular}{lr}
\toprule
Check category&Count\\
\midrule
Independent rational coefficient comparisons, weights $0$--$32$&8,350\\
Exact support and sign checks on positive highest indices&8,317\\
Full identities $A_m=N_m+\D H_m$, weights $0$--$28$&29\\
Euler-homotopy versus recursive primitive, weights $0$--$18$&19\\
Appell derivative identities, weights $1$--$32$&32\\
Mixed-square finite-sum comparisons&1,684\\
Deliberately corrupted certificates rejected&2\\
Exact analytic endpoint-algebra assertions&20\\
\bottomrule
\end{tabular}
\end{center}
The two corrupted certificates change the displayed weight-eleven and
cubic weight-fourteen coefficients while retaining their independently
computed primitives; the polynomial identity fails as it should.
The exact endpoint assertions include Bernoulli-polynomial evaluations
at thirteen nonpositive spectral orders, the residue, and the rational
and harmonic coefficient algebra of the worked examples.
The supplied CSV contains all 8,350 coefficient slots, including zeros.
The JSON certificates contain complete $A_m,N_m,H_m$ data at the five
selected weights, not just numerical evaluations.

\subsection{Numerical diagnostics and their limits}
The analytic script was run with 65 decimal digits using Python 3.13.5,
mpmath 1.3.0, and SymPy 1.14.0. All 20 diagnostics passed their declared
local tolerances. They include real and complex Hurwitz orders,
trigamma, $\gamma_0$, $\gamma_1$, a log-sine integral, a finite
polylogarithmic endpoint formula, mixed beta moments, four logarithmic
orders, an anchored Gamma primitive, and an ordinary Gamma endpoint
identity.

For independent quadrature values $L,R$, the recorded residual is
$|L-R|/(1+|L|+|R|)$. The largest recorded residual is about
$4.82\times10^{-34}$, in the refinement comparison between 32- and
48-point Gauss--Legendre quadrature for the $\gamma_1$ integral; its
declared tolerance is $10^{-31}$. The finer quadrature versus the
finite endpoint value differs by about $2.68\times10^{-52}$ in the
same scaled measure. These are floating-point diagnostics, not rigorous
interval enclosures, and 65-digit working precision does not mean that
every integral has 65 certified digits.

For orientation, the four ordinary integrals in
\eqref{ubg:eq:trigamma-example}, \eqref{ubg:eq:gamma0-example},
\eqref{ubg:eq:gamma1-example}, and \eqref{ubg:eq:logsine-example} begin
\[
 \begin{gathered}
 0.83786519860377187790,\qquad -0.47562629189112816162,\\
 -0.11146849703321255941,\qquad 0.11670482124545205075,
 \end{gathered}
\]
respectively. The finite identities and their proofs, rather than these
decimal values, are the research statements.

\paragraph{Replay and integration.}
From the package root, run
\begin{verbatim}
python verification/check_exact.py
python verification/check_analytic.py
pdflatex -interaction=nonstopmode -halt-on-error article.tex
pdflatex -interaction=nonstopmode -halt-on-error article.tex
\end{verbatim}
The exact coefficient checks need only the Python standard library.
The analytic checks require the two packages listed above.
The archive includes an additive integration plan and a compact theorem
fragment with unique \src{ubg:} labels. No connected repository files
were modified. The result has not been formalized in a proof assistant.

\section{Further research questions after the coefficient theorem}
\label{ubg:sec:questions}

The all-monomial theorem changes the frontier: finding the next mixed
Gamma coefficient is no longer an open problem for this hierarchy.
The following targets concern genuinely additional structure rather
than an already solved sector.

\paragraph{1. Exact denominators and arithmetic normalization.}
Determine the least positive integer $D_m$ for which $D_mN_m$ has
integer coefficients, and characterize the denominator of an individual
coefficient from its exponent vector. The first cubic example
$-3185/3$ rules out a blanket integrality statement. The positive finite
sum \eqref{ubg:eq:positive-coefficient} gives a concrete starting point
for prime-adic cancellation identities. A satisfactory answer should
specify exact valuations or divisibility laws, not merely denominator
bounds.

\paragraph{2. Analytic summation of the full normal-form series.}
For an analytic function $u_0(x)$, determine when the formal infinite
sum underlying \eqref{ubg:eq:generator} converges after all jet
substitutions, and identify its analytic value. The theorem proves
entire dependence for every \emph{fixed finite jet}, not convergence
of the infinitely many jets of $\psi(x)$. It would be particularly
useful to obtain an exact integral representation on a specified
$(x,t)$ domain and the corresponding summed endpoint primitive.

\paragraph{3. General centered profiles and products of Gamma ratios.}
Replace $K_t(\D)$ by an even profile with leading term $t$, or consider
finite products of centered Gamma ratios. The affine predecessor
already supplies square symbols for general symmetric profiles
\cite{Affine}. Differential superposition still provides a route to
polynomial backgrounds. Determine which profiles admit an equally
simple full-coefficient measure, and give necessary and sufficient
conditions for the exact sign law to survive. Positivity of the Gamma
beta kernel is a special input, not a formal consequence for every
profile.

\paragraph{4. Arithmetic selection by coefficient weights.}
The predecessor's Bernoulli selectors diagonalize certain polynomial
Hurwitz integrals into single shifted zeta terms \cite{Affine}.
Classify the universal weights $P_M$ that lie in one selector line,
or in a prescribed low-dimensional selector span. This asks which
actual Gamma coefficient patterns produce a small number of zeta
values, rather than redoing the existing classification of arbitrary
even polynomial weights.

\paragraph{5. Exact values of individual nonlinear polygamma integrals.}
The universal normal form supplies all linear relations arising from
this Gamma hierarchy, but it does not evaluate every surviving
functional integral. A concrete target is to identify the natural
function algebra for
$\int_1^2\psi'(x)(\psi'''(x))^3\dd x$, or for a parameterized version
with independently shifted arguments. Investigate a finite reduction
to multiple Hurwitz or multiple-polylogarithmic quantities with fixed
branches, and distinguish an actual evaluation from a relation in
which new unevaluated integrals remain.

\paragraph{6. Endpoint collisions in the Stieltjes and Lerch transforms.}
The present formulas assume $A>B>0$. Continue the exact endpoint
identities as $A\downarrow B$, keeping a specified coordinate and
subtraction convention. Which weights kill every local singular term,
and which leave a nonzero contact contribution? The ordinary log-sine
limit at order one does not settle the distributional limits at
nonpositive polylogarithm orders. The repository's boundary-contact
work should be used rather than silently identifying the two regimes.

\paragraph{7. A proof-carrying compiler.}
Formalize the terminating functional decomposition, weighted quadratic
symbol, differential superposition, and highest-variation extraction.
The theorem splits naturally into finite polynomial identities and
coefficientwise formal-series arguments. A first practical milestone
is a proof-assistant checker for the supplied rational certificates;
a separate milestone is a formal all-order proof. Certified interval
versions of the analytic diagnostics would add numerical evidence but
would not substitute for either formal theorem.

\paragraph{8. Preserve the original period and tail-rank targets.}
The exact Gaussian $S_6$ and revised $S_8$ candidates still require
branch-controlled polylogarithmic identities with their complete
lower-weight constants. The general unequal harmonic-tail moment-prefix
rank question still requires a genuine determinant or independence
argument. The present coefficient theorem gives additional exact
identities and endpoint tools, but no implication resolving those
separate questions is claimed.

\subsection*{Conclusion}
The normal-form coefficient problem has a uniform answer: its highest
quadratic variation on a polynomial background is a compact exponential
beta integral, and two integrations in the highest jet recover every
multiplicity. This yields a finite rational formula for every functional
coefficient, a complete support criterion, and a uniform sign law.
Keeping the primitive gives exact Gamma--harmonic endpoint identities;
retaining the polynomial weight gives finite zeta, Stieltjes, and
polylogarithmic transforms. The distinctions between formal reduction,
ordinary integration, joint meromorphic continuation, and numerical
verification remain explicit throughout.

\begin{thebibliography}{99}
\raggedright
\bibitem{Manuscript}
ProveIt Contributors, \emph{Polylogarithms and their Arithmetic Bridges}.
Canonical manuscript, especially the normal-form material and chapters
\src{07-ibp-normal-forms.tex} and
\src{07-differential-polynomial-moments.tex}.
Targeted repository reads, 11 October 2026.
\url{https://github.com/VladimirReshetnikov/ProveIt/tree/main/Analysis/Polylogarithms/docs/manuscript}.
The incoming-directory inspection was at
\url{https://github.com/VladimirReshetnikov/ProveIt/tree/main/docs/incoming}.

\bibitem{HTG}
\emph{Exact Harmonic-Tail Reconstruction and Variational Gamma Identities}.
ProveIt working report, 11 October 2026. Especially
\src{sections/04_gamma_sectors.tex}; article blob identity
\src{445e67f0cd60313113387d52d44d4ddc6eeaf244}.
\url{https://github.com/VladimirReshetnikov/ProveIt/tree/main/Analysis/Polylogarithms/docs/reports/stieltjes-harmonic-resonance/continuations/harmonic-tail-gamma}.

\bibitem{Affine}
\emph{Affine Polarization of Gamma Normal Forms: All derivative-square
coefficients, polylogarithmic kernels, and Stieltjes averages}.
Research continuation prepared for Vladimir Reshetnikov, 11 October 2026.
Full TeX source retrieved from the user's file library as
\src{ProveIt_Affine_Gamma_Polarization_2026-10-11.tex}.
Sections 2--3 and 11 are the direct normal-form input and open-question
source. The exact source fingerprint and inspection limits appear in
Section~\ref{ubg:sec:audit} and the package manifest.

\bibitem{GRR}
Li Guo, Georg Regensburger, and Markus Rosenkranz,
\emph{On integro-differential algebras}.
Journal of Pure and Applied Algebra \textbf{218} (2014), 456--471.
\url{https://doi.org/10.1016/j.jpaa.2013.06.015}.
Author preprint and journal-reference record:
\url{https://arxiv.org/abs/1212.0266}.
Section 4.4.1 supplies the classical functional-monomial construction.

\bibitem{EM1}
Olivier R. Espinosa and Victor H. Moll,
\emph{On some definite integrals involving the Hurwitz zeta function}.
Author preprint, 2000.
\url{https://arxiv.org/abs/math/0012078}.
Cited for the established polynomial-weight Hurwitz, Bernoulli,
and log-Gamma integration programme.

\bibitem{EM2}
Olivier R. Espinosa and Victor H. Moll,
\emph{On some integrals involving the Hurwitz zeta function: part 2}.
Author preprint, 2001.
\url{https://arxiv.org/abs/math/0107082}.
Cited for established antiderivative and generalized-polygamma methods.

\bibitem{DLMFpsi}
NIST Digital Library of Mathematical Functions, \S5.15,
\emph{Polygamma Functions}. Consulted 11 October 2026.
\url{https://dlmf.nist.gov/5.15}.

\bibitem{DLMFzeta}
NIST Digital Library of Mathematical Functions, \S25.11,
\emph{Hurwitz Zeta Function}. Consulted 11 October 2026.
\url{https://dlmf.nist.gov/25.11}.
Definitions, parameter derivatives, integer shifts, Bernoulli values,
Mellin representation, and log-Gamma specialization.

\bibitem{DLMFintegers}
NIST Digital Library of Mathematical Functions, \S25.6,
\emph{Integer Arguments}. Consulted 11 October 2026.
\url{https://dlmf.nist.gov/25.6}.

\bibitem{DLMFlerch}
NIST Digital Library of Mathematical Functions, \S25.14,
\emph{Lerch's Transcendent}. Consulted 11 October 2026.
\url{https://dlmf.nist.gov/25.14}.

\bibitem{DLMFpolylog}
NIST Digital Library of Mathematical Functions, \S25.12,
\emph{Polylogarithms}. Consulted 11 October 2026.
\url{https://dlmf.nist.gov/25.12}.
\end{thebibliography}
\end{document}
